\documentclass[10pt,a4paper]{amsart}
\usepackage[margin=19mm]{geometry}
\usepackage{amsmath,amssymb,mathtools}
\usepackage[scr=rsfs,cal=euler]{mathalfa}
\usepackage{xfrac}
\usepackage[unicode,hidelinks,bookmarks=false]{hyperref}
\newcommand{\ie}{i.e.\ }
\newcommand{\cf}{cf.\ }
\newcommand{\resp}{respectively\ }
\newcommand{\Subsection}{\subsection}
\providecommand{\nobreakdash}{\nobreak\hbox{-}}
\setlength{\emergencystretch}{2em}
\multlinegap0pt
\usepackage{amssymb}
\usepackage{mathtools}
\usepackage{bm}
\usepackage[shortlabels]{enumitem}
\usepackage{xcolor}
\newtheorem{proposition}{Proposition}
\newcommand{\C}{\mathbb C}
\newcommand{\E}{\mathbb E}
\newcommand{\Herm}{\mathrm{Herm}}
\DeclareMathOperator{\Per}{Per}
\newcommand{\Pp}{\mathbb P}
\newcommand{\R}{\mathbb R}
\DeclareMathOperator{\Tr}{Tr}
\newcommand{\dd}{\mathrm d}
\DeclareMathOperator{\diag}{diag}
\newcommand{\ii}{\mathrm i}
\title{Rotated semicircle laws for permanental roots of Gaussian random matrices}
\author{Renjie Feng, Dong Yao}
\date{Source: arXiv:2608.26038v1 (CC BY 4.0)}
\begin{document}
\maketitle

\noindent\emph{Source and licence.} Rotated semicircle laws for permanental roots of Gaussian random matrices. Version arXiv:2608.26038v1 (CC BY 4.0), distributed by arXiv under the Creative Commons Attribution 4.0 International licence, http://creativecommons.org/licenses/by/4.0/, as recorded in the arXiv licence metadata for this exact version and shown on the abstract page https://arxiv.org/abs/2608.26038v1. Full licence notice: \texttt{licenses/arxiv-2608.26038-CC-BY-4.0.txt}.\par\smallskip

\noindent\textit{Editorial scope.} Counted unit: the article's proposition prop:scalar-permanent-duality (source0.tex lines 489--508). The article's complete original proof of it is delivered with it (source0.tex lines 510--627, one \texttt{proof} environment of 118 lines). No mathematics is written, repaired or re-derived: the statement and the proof are the article's own lines, in the article's own notation, and no retained line is substituted or re-flowed.  Retained uncounted as the source-local context the proof needs: lines 360--365; lines 371--375; lines 392--397; lines 404--410.  Nothing was omitted from any retained span, so the package carries no omission record.  No retained line contains a citation, so the wrapper carries no bibliography and no bibliography entry.  No retained line contains a figure environment, an image-inclusion command or drawing code, so no image is delivered and the package declares no asset dependency.  Editorial formatting only, in addition: an AMS \texttt{amsart} preamble that restates the article's own declarations and macros used by the retained lines, namely the environments \texttt{proposition} on separate counters, together with the standard packages those lines need.  The retained environments print in this document as Proposition 1 in order of appearance, whereas the article prints the same items with the numbering of its own full document.  The complete native source of the article as distributed in the arXiv e-print for this exact version is bundled unchanged as \texttt{sources/208614/source0.tex}.
\begin{equation}\label{eq:aux-q-law}
  \dd\gamma_{q,N}(q)
  :=
  \frac{N^2}{2\pi^2}
  \exp\!\left(-\frac N2\Tr q^2\right)\dd q.
\end{equation}

\begin{equation}\label{eq:eta-law}
  \dd\gamma_{\eta,N}(\eta)
  :=
  \frac N\pi e^{-N|\eta|^2}\dd^2\eta.
\end{equation}

\begin{equation}\label{eq:GUE-per-duality}
  \E[P_N(\mu_1)P_N(\mu_2)]
  =
  \int_{\Herm_2(\C)}
  [\Per(M-q)]^N\,\dd\gamma_{q,N}(q).
\end{equation}

\begin{equation}\label{eq:GOE-per-duality}
  \E[P_N(\mu_1)P_N(\mu_2)]
  =
  \int_{\Herm_2(\C)}\int_{\C}
  \bigl(\Per(M-q)+|\eta|^2\bigr)^N
  \,\dd\gamma_{\eta,N}(\eta)\,\dd\gamma_{q,N}(q).
\end{equation}

\begin{proposition}\label{prop:scalar-permanent-duality}
For every \(x\in\R\),
\begin{align}
  H_N\sim\Pp_{2,N}:\qquad
  &\E[P_N(\ii x)P_N(-\ii x)]
  =
  \int_{\R^2}
  \prod_{\ell=1}^2(x-\ii y_\ell)^N
  \,\dd\Gamma_{2,N}(y_1,y_2),
  \label{eq:GUE-per-scalar}\\
  H_N\sim\Pp_{1,N}:\qquad
  &\E[P_N(\ii x)P_N(-\ii x)]
  =
  \int_{\R^2}
  \prod_{\ell=1}^2(x-\ii y_\ell)^N
  \,\dd\Gamma_{4,N}(y_1,y_2).
  \label{eq:GOE-per-scalar}
\end{align}
The integrals are absolutely convergent.
\end{proposition}

\begin{proof}
We first consider the GUE.  In \eqref{eq:GUE-per-duality}, take
\[
  M=\diag(\ii x,-\ii x).
\]
Make the change of variable \(q_{11}\mapsto-q_{11}\), which preserves the
measure \(\gamma_{q,N}\), and then relabel the reflected coordinate as
\(q_{11}\).  A direct calculation gives
\[
  \Per(M-q)
  =
  x^2-\ii x(q_{11}+q_{22})-q_{11}q_{22}+|q_{12}|^2.
\]
Write
\[
  s=\frac{q_{11}+q_{22}}2,
  \qquad
  v=\frac{q_{11}-q_{22}}2,
  \qquad
  u=\Re q_{12},
  \qquad
  \omega=\Im q_{12},
\]
and put
\[
  r=(v^2+u^2+\omega^2)^{1/2}.
\]
Then
\[
  \Per(M-q)
  =
  (x-\ii(s+r))(x-\ii(s-r)),
  \qquad
  \Tr q^2=2(s^2+r^2).
\]
The absolute Jacobian of the linear transformation
\((s,v)\mapsto(q_{11},q_{22})\) is \(2\).  Hence
\eqref{eq:aux-q-law} becomes
\[
  \frac{N^2}{\pi^2}e^{-N(s^2+r^2)}
  \,\dd s\,\dd v\,\dd u\,\dd\omega.
\]
Using polar coordinates in \(\R^3\), together with \(|S^2|=4\pi\), gives
\begin{equation}\label{eq:GUE-s-r-density}
  \frac{4N^2}{\pi}r^2e^{-N(s^2+r^2)}\,\dd s\,\dd r,
  \qquad s\in\R,\quad r\ge0.
\end{equation}
Now set
\[
  y_1=s+r,
  \qquad
  y_2=s-r.
\]
This maps \(\R\times[0,\infty)\) onto the half-plane
\(\{(y_1,y_2)\in\R^2:y_1\ge y_2\}\), up to boundary sets of Lebesgue
measure zero.  Moreover,
\[
  \dd s\,\dd r=\frac12\,\dd y_1\,\dd y_2,
  \qquad
  r=\frac{y_1-y_2}{2}.
\]
Therefore the push-forward of \eqref{eq:GUE-s-r-density} has density
\begin{equation}\label{eq:ordered-beta2-density}
  \frac{N^2}{2\pi}(y_1-y_2)^2
  e^{-\frac N2(y_1^2+y_2^2)}\,\dd y_1\,\dd y_2
\end{equation}
on \(y_1\ge y_2\).  The function
\[
  (y_1,y_2)\longmapsto
  (x-\ii y_1)^N(x-\ii y_2)^N
\]
is symmetric in \(y_1,y_2\).  Since \(Z_{2,N}=4\pi/N^2\), integrating
this symmetric function against \eqref{eq:ordered-beta2-density} on
\(y_1\ge y_2\) is equal to integrating it against \(\Gamma_{2,N}\) on
all of \(\R^2\).  This proves \eqref{eq:GUE-per-scalar}.

We next consider the GOE.  Starting from \eqref{eq:GOE-per-duality}, make
the same reflection \(q_{11}\mapsto-q_{11}\), and write
\[
  \eta=\eta_1+\ii\eta_2.
\]
With \(s,v,u,\omega\) as above, set
\[
  r=
  \bigl(v^2+u^2+\omega^2+\eta_1^2+\eta_2^2\bigr)^{1/2}.
\]
Then
\[
  \Per(M-q)+|\eta|^2
  =
  (x-\ii(s+r))(x-\ii(s-r)).
\]
Combining \eqref{eq:aux-q-law} and \eqref{eq:eta-law}, and including the
Jacobian \(2\) from \((s,v)\mapsto(q_{11},q_{22})\), gives
\[
  \frac{N^3}{\pi^3}e^{-N(s^2+r^2)}
  \,\dd s\,\dd v\,\dd u\,\dd\omega\,\dd\eta_1\,\dd\eta_2.
\]
Using polar coordinates in \(\R^5\), together with
\(|S^4|=8\pi^2/3\), gives
\begin{equation}\label{eq:GOE-s-r-density}
  \frac{8N^3}{3\pi}r^4e^{-N(s^2+r^2)}\,\dd s\,\dd r,
  \qquad s\in\R,\quad r\ge0.
\end{equation}
Under the same transformation \(y_1=s+r\), \(y_2=s-r\), the push-forward
of \eqref{eq:GOE-s-r-density} has density
\begin{equation}\label{eq:ordered-beta4-density}
  \frac{N^3}{12\pi}(y_1-y_2)^4
  e^{-\frac N2(y_1^2+y_2^2)}\,\dd y_1\,\dd y_2
\end{equation}
on the half-plane \(y_1\ge y_2\).  Since \(Z_{4,N}=24\pi/N^3\) and the
integrand is symmetric in \(y_1,y_2\), integration against
\eqref{eq:ordered-beta4-density} on \(y_1\ge y_2\) equals integration
against \(\Gamma_{4,N}\) on all of \(\R^2\).  This proves
\eqref{eq:GOE-per-scalar}.

 
\end{proof}

\end{document}
